\documentclass[12pt]{article}
\usepackage{lmodern}
\usepackage[utf8]{inputenc}
\usepackage{amsmath, amssymb}
\usepackage{tabularx}
\usepackage{geometry}
\usepackage{multirow}
\usepackage{graphicx}
\usepackage[table]{xcolor}
\usepackage{draftwatermark}

\title{Linearly Independent in Different Sizes of Matrix}

\SetWatermarkText{WongAndyStarSB}
\SetWatermarkScale{0.6}
\SetWatermarkColor{blue!03}

\geometry{margin=1in}

% centralized style management
\newcommand{\myvec}[1]{\vec{#1}}
\definecolor{mycolorN} {HTML}{007514}   % hsl: 130, 100, 23
\definecolor{mycolorNr}{HTML}{007575}   % hsl: 180, 100, 23
\definecolor{mycolorNc}{HTML}{4e7500}   % hsl:  80, 100, 23
\renewcommand{\arraystretch}{2}

\begin{document}

\section*{Linearly Independent in Different Sizes of Matrix}
\raggedright

for a matrix $A$ in $\textcolor{mycolorNr}{n_r} \times \textcolor{mycolorNc}{n_c}$, \newline
if the \emph{column vectors} of A are {\large \emph{linearly independent}}, then

\medskip


\begin{tabularx}{\textwidth}{XXc}
    \arrayrulecolor{black!30} \hline \arrayrulecolor{black}
        $\textcolor{mycolorN}{n} = \textcolor{mycolorNr}{n_r} = \textcolor{mycolorNc}{n_c}$ \textbf{(Square)} 
        & \textbf{$\textcolor{mycolorNr}{n_r} > \textcolor{mycolorNc}{n_c}$ (Tall)} 
        & \textbf{$\textcolor{mycolorNr}{n_r} < \textcolor{mycolorNc}{n_c}$ (Fat)}
    \\ \hline
        $\iff$ only trivial solution \newline
            \phantom{$\iff$} in $A \myvec{x} = \myvec{0}$ 
        & $\iff$ only trivial solution \newline
            \phantom{$\iff$} in $A\myvec{x} = \myvec{0}$ 
        & \multirow[c]{9.5}{*}{\textbf{NOT POSSIBLE}}
    \\ \arrayrulecolor{black!30} \cline{1-2} \arrayrulecolor{black}
        $\iff Col A = \mathbb{R}^{\textcolor{mycolorN}{n}}$ \newline 
            \phantom{$\iff$} $Nul A = \{\myvec{0}\}$ \newline 
            \phantom{$\iff$} $rank A = \textcolor{mycolorN}{n}$ (full) 
        & $\iff Col A = \mathbb{R}^{\textcolor{mycolorNr}{n_r}}$ \newline 
            \phantom{$\iff$} $Nul A = \{\myvec{0}\}$ \newline 
            \phantom{$\iff$} $rank A = \textcolor{mycolorNc}{n_c}$ (full) 
        & 
    \\ \arrayrulecolor{black!30} \cline{1-2} \arrayrulecolor{black}
        $\iff \myvec{x}$ and $\myvec{v}$ are \newline
            \phantom{$\iff$} \emph{one-to-one} \newline
            \phantom{$\iff$} in $A\myvec{x} = \myvec{v}$
        & $\iff \myvec{x}$ and $\myvec{v}$ are \newline
            \phantom{$\iff$} \emph{one-to-one} \newline
            \phantom{$\iff$} in $A\myvec{x} = \myvec{v}$
        &
    \\ \arrayrulecolor{black!70} \cline{1-2} \arrayrulecolor{black}
        $\iff$ unique solution \newline
            \phantom{$\iff$} in $A\myvec{x} = \myvec{v}$ 
        & $\iff$ 0 or 1 solution \newline
            \phantom{$\iff$} in $A\myvec{x} = \myvec{v}$ \newline 
            \phantom{$\iff$} (may not exist, \textsuperscript{1} \newline
            \phantom{$\iff$(} if exist then unique \textsuperscript{2}) 
        & 
    \\ \arrayrulecolor{black!30} \cline{1-2} \arrayrulecolor{black}
        $\iff \det A \neq 0$ \newline 
            \phantom{$\iff$} \emph{Invertible} \newline 
            \phantom{$\iff$} $A^{-1}$
        & \emph{(not applicable)} 
        & 
    \\ \hline
\end{tabularx}


\vfill


\begin{minipage}[t]{0.48\textwidth}
    \raggedright
    
    \textsuperscript{1}/ \\
    \textbf{Consider} $A\myvec{x}=\myvec{v}$: \\
    it represents $\textcolor{mycolorNr}{n_r}$ equations with \\
    \phantom{it represents }$\textcolor{mycolorNc}{n_c}$ unknowns\\
    in other words:\\
    \hspace*{1em} num of restrictions $>$ num of variables.\\
    \hspace*{1em} $\rightarrow$ no solutions \\
    \hspace*{1em} \phantom{$\rightarrow$} \small{(unless the restrictions overlaps)}
\end{minipage}
\begin{minipage}[t]{0.48\textwidth}
    \raggedright

    \textsuperscript{2}/ \\
    \textbf{Consider} the case where\\ 
    $A\myvec{x} = \myvec{v}$ is consistent\\
    Let $\myvec{x}_1$ and $\myvec{x}_2$ be solutions of $\myvec{x}$
    \begin{align*}
        A\myvec{x}_1 = \myvec{v} \text{ and } &A\myvec{x}_2 = \myvec{v}\\
        A\myvec{x}_1 - A\myvec{x}_2 &= \myvec{v} - \myvec{v}\\
        A(\myvec{x}_1 - \myvec{x}_2) &= \myvec{0}\\
        \myvec{x}_1 - \myvec{x}_2 &= \myvec{0} && (Nul A = \{\myvec{0}\})\\
        \myvec{x}_1 &= \myvec{x}_2
    \end{align*}
    $\rightarrow$ $A\myvec{x} = \myvec{v}$ has unique solution for $\myvec{x}$ if \\
    \phantom{$\rightarrow$} the equation is consistent
\end{minipage}

\end{document}